% =============================================================
%  chapter1_abstract.tex
%  Shared abstract text for Chapter 1 / TubeMethodPaper.tex.
%
%  Kept separate from the surrounding \begin{abstract}...\end{abstract}
%  wrapper because that environment is NOT class-agnostic: under
%  copernicus.cls it renders in place as the paper's abstract, but
%  under caltech_thesis.cls the `abstract` environment is redefined
%  to break out into a new unnumbered "Abstract" chapter (via
%  \extrachapter), which is only appropriate once, at the front of
%  the whole dissertation (see dissertation.tex's own Abstract
%  chapter) -- not again in the middle of Chapter 1. So each wrapper
%  supplies its own environment/heading around this same raw text.
%
%  Edit this file (not TubeMethodPaper.tex or chapter1.tex directly)
%  to keep the journal paper and the dissertation chapter in sync.
% =============================================================
Helium isotope studies and U-Th/He geochronology of sulfur-bearing and hydrous minerals offer direct windows into the timing and context of hydrothermal mineralization, provided the conditions governing radiogenic helium retention and loss in nature are known. However, assessing helium retentivity at elevated temperatures is currently hampered by these minerals\textquotesingle{} instability under the high-temperature, high-vacuum conditions of conventional helium step-heating experiments. This paper introduces a novel method for conducting helium step-heating experiments under hydrothermal and sulfur-saturated conditions, designed to extend our knowledge of helium diffusion kinetics in hydrous and sulfur-bearing mineral phases. In this approach, silica glass ampoules encapsulate the mineral of interest, maintaining internally sulfur- or water vapor-saturated conditions against the ultra-high vacuum of the surrounding extraction line, while the high helium permeability of fused silica permits quantification of evolved helium using standard temperature control and mass spectrometry procedures. This work presents the experimental framework and key equations necessary to deconvolve the processes of helium permeation through the silica ampoule from the thermally activated degassing of helium from the encapsulated mineral. Permeation lag is quantified with a first-order permeation model and removed by one of two routes: the \emph{permeation step method}, which measures the ampoule's residual helium directly with low-temperature permeation steps, or, where those yields fall to blank level, the \emph{bounding method}, which constrains that residual from above by assuming steady-state permeation during long steps. Both rest on the ampoule's permeation rate constant, which is measured during the experiment itself rather than taken from the literature. Validation experiments measuring helium permeation through empty ampoules and ampoules containing only water, sulfur, or Durango apatite demonstrate that permeation lag is readily correctable in well-designed experiments and characterize the extent to which water and sulfur influence the helium permeability of fused silica. While sulfur-saturated experiments are straightforward to carry out with this method, several tractable complications are identified for experiments when significant water is present. This method expands standard experimental capabilities to encompass mineral phases that decompose under high vacuum, including sulfides and iron oxyhydroxides, opening new avenues for thermochronometric investigation of hydrothermal systems.
